\title{Groups Are Factors: Learning Linear PDEs\\by Clustering Modes on the Factors of the Symbol}
\author{Flavio Pricoco\thanks{\email{flaflapr@gmail.com}}}
\headers{Groups are factors: learning linear PDEs}{F. Pricoco}

\begin{document}
\maketitle

\begin{abstract}
Methods that learn partial differential equations from data usually regress a time derivative on a library of terms.
We argue that for linear equations with constant coefficients the natural unit is not the term but the factor of the
symbol. Every Fourier mode of a solution lies on the zero set of one factor, so learning the equation means grouping
the modes by factor. We prove that every solution splits into one component per factor, uniquely up to an overlap
counted by the intersection points of the zero sets; that additive terms admit no such split; that the split is stable
exactly when the factors' solution spaces are well conditioned (for two factors, when their Friedrichs angle is
positive); and that finite data can support but never certify an equation. The resulting algorithm groups twice:
within one sample of the field it clusters Fourier modes onto the curves of the factors, and across posterior samples
it clusters the factors with a Dirichlet-process mixture. Its output is a distribution over equations with intervals,
the separation of each group, and a label stating what the data support. On 285 random equations never used in
development it returns the right or an equivalent equation on 96.5\% and is never wrong without warning, while weak
SINDy is right on 28.4\% and wrong without warning on 204. On a ten-equation development benchmark it recovers all
thirty factor sets, including the products of factors on which the best library regression fails.
\end{abstract}

\begin{keywords}
equation discovery, operator factorization, dispersion relation, Friedrichs angle, matrix pencil, Dirichlet process,
uncertainty quantification
\end{keywords}

\begin{MSCcodes}
35E20, 62F15, 65M32, 68T05
\end{MSCcodes}

\section{Introduction}\label{sec:intro}
Consider a quantity that is carried by a flow and diffuses at the same time, observed with noise. A natural model is
\begin{equation}\label{eq:intro-tpd}
  (\partial_t+c\,\partial_x)(\partial_t-\nu\,\partial_x^2)\,u=0,
\end{equation}
whose solutions are sums of a travelling profile and a diffusing one. A method that regresses $u_t$ on a library of
terms such as $u_x$, $u_{xx}$, $u_{xxx}$ looks for one equation $u_t=\sum_k\beta_k\,\partial_x^ku$. The data obey the
\emph{product} of two such equations, and no single sum of terms fits them: every library-regression method we tested
returns a wrong equation for \eqref{eq:intro-tpd} at $5\%$ noise, and none of them warns.

\paragraph{The idea} In Fourier variables a constant-coefficient operator is a polynomial in $(\tau,\xi)$, with
$\tau\leftrightarrow\partial_t$ and $\xi\leftrightarrow\partial_x$; the operator in \eqref{eq:intro-tpd} is
$(\tau+c\,\xi)(\tau-\nu\,\xi^2)$. A Fourier mode $e^{\lambda t+i\kappa x}$ solves the equation exactly when $(\lambda,i\kappa)$
is a zero of this polynomial, that is, a zero of one of its factors. So the modes of every solution lie on a few curves,
one per factor: here the line $\lambda=-ic\kappa$ (transport) and the parabola $\lambda=-\nu\kappa^2$ (diffusion), as in
\cref{fig:branches}. \emph{Learning the equation means finding these curves, which is a clustering problem whose
clusters are curves rather than points.} This paper develops that idea from the algebra to an algorithm with calibrated
uncertainty.

\begin{figure}[tbp]
\centering
\begin{tikzpicture}[x=0.40cm,y=0.40cm,>=Stealth,font=\footnotesize]
  \begin{scope}
    \draw[->,gray] (-4.5,0)--(4.8,0) node[right,black]{$\kappa$};
    \draw[->,gray] (0,-4.5)--(0,4.8) node[above,black]{$\operatorname{Im}\lambda$};
    \draw[very thick,color={rgb,255:red,42;green,120;blue,214}] (-4.2,4.2)--(4.2,-4.2);
    \draw[very thick,color={rgb,255:red,235;green,104;blue,52},dashed] (-4.2,0)--(4.2,0);
    \foreach \k in {-3,-2,-1,1,2,3} \fill[color={rgb,255:red,42;green,120;blue,214}] (\k,-\k) circle (2.6pt);
    \foreach \k in {-2,-1,1,2} \fill[color={rgb,255:red,235;green,104;blue,52}] (\k,0) +(-2.4pt,-2.4pt) rectangle +(2.4pt,2.4pt);
    \node[anchor=west,align=left] at (1.2,3.3) {transport\\$\lambda=-ic\kappa$};
    \node at (0,-5.8) {(a) oscillation};
  \end{scope}
  \begin{scope}[xshift=10.6cm]
    \draw[->,gray] (-4.5,0)--(4.8,0) node[right,black]{$\kappa$};
    \draw[->,gray] (0,-4.5)--(0,1.6) node[above,black]{$\operatorname{Re}\lambda$};
    \draw[very thick,color={rgb,255:red,42;green,120;blue,214},dashed] (-4.2,0)--(4.2,0);
    \draw[very thick,color={rgb,255:red,235;green,104;blue,52},domain=-3.8:3.8,samples=60] plot (\x,{-0.3*\x*\x});
    \foreach \k in {-3,-2,-1,1,2,3} \fill[color={rgb,255:red,42;green,120;blue,214}] (\k,0) circle (2.6pt);
    \foreach \k in {-2,-1,1,2} \fill[color={rgb,255:red,235;green,104;blue,52}] (\k,{-0.3*\k*\k}) +(-2.4pt,-2.4pt) rectangle +(2.4pt,2.4pt);
    \node[anchor=west,align=left] at (1.8,1.6) {diffusion\\$\lambda=-\nu\kappa^2$};
    \node at (0,-5.8) {(b) damping};
  \end{scope}
\end{tikzpicture}
\caption{The idea of the paper on $(\partial_t+c\,\partial_x)(\partial_t-\nu\,\partial_x^2)u=0$. Every Fourier mode
$e^{\lambda t+i\kappa x}$ of a solution lies on the zero set of one factor of the symbol $(\tau+c\xi)(\tau-\nu\xi^2)$:
circles on the line of the transport factor, squares on the parabola of the diffusion factor (schematic;
\cref{fig:directions} shows the modes found in noisy data). Learning the equation means learning these two
groups.}\label{fig:branches}
\end{figure}

\paragraph{What the theory says} Three questions decide whether this clustering is well posed, and
\cref{sec:theory} answers each. \emph{Are the groups really the factors?} Yes: every solution splits into one
component per factor, uniquely up to an overlap whose dimension is the number of intersection points of the zero sets
(\cref{thm:factor}); additive terms, by contrast, admit no such split (\cref{prop:terms}). \emph{How well can the
groups be separated?} Exactly as well as the Friedrichs angle between the factors' solution spaces allows
(\cref{thm:hm}). \emph{What can finite data establish?} They can support a factored equation but never certify it
(\cref{prop:certify}); this is why the algorithm reports labels, not ``identified''.

\paragraph{What the algorithm does} It groups twice (\cref{fig:roadmap}). \emph{Within one sample} of the field it
reads the temporal rates of every wavenumber (the points of \cref{fig:branches}) and clusters them onto the curves of
candidate factors. \emph{Across posterior samples} it clusters the factors found in each sample, with a Dirichlet-process
mixture, so that the uncertainty of the data becomes a distribution over equations. Every decision is finally checked
against the raw observations, and the output carries intervals, the separation of each group and a label.

\begin{figure}[tbp]
\centering
\begin{tikzpicture}[
  node distance=3.5mm and 4.5mm,
  box/.style={draw=black!55, rounded corners=2pt, align=center, font=\footnotesize, inner sep=3pt,
              minimum height=10mm, text width=20mm, fill=white},
  grp/.style={box, fill=blue!8, very thick},
  thy/.style={font=\scriptsize\itshape, text=black!60, align=center},
  arr/.style={-{Stealth[length=2mm]}, black!60, thick}]
\node[box] (data) {noisy data\\$\{(t_i,x_i,y_i)\}$};
\node[box, right=of data] (gp) {field samples\\$u^{(1)},\dots,u^{(\Nsamp)}$};
\node[box, right=of gp] (modes) {modes $(\kappa,\lambda)$\\per wavenumber};
\node[grp, right=of modes] (fac) {\textbf{grouping I}\\modes $\to$ factors};
\node[grp, below=8mm of fac] (cloud) {\textbf{grouping II}\\factors over samples};
\node[box, left=of cloud] (clus) {distribution\\over equations};
\node[box, left=of clus] (split) {components $u_j$,\\separation $\gN$};
\node[box, left=of split] (out) {equation with\\intervals and label};
\draw[arr] (data)--(gp); \draw[arr] (gp)--(modes); \draw[arr] (modes)--(fac);
\draw[arr] (fac)--(cloud); \draw[arr] (cloud)--(clus); \draw[arr] (clus)--(split); \draw[arr] (split)--(out);
\node[thy, above=1mm of gp] {\S\ref{sec:modes}};
\node[thy, above=1mm of modes] {\S\ref{sec:modes}};
\node[thy, above=1mm of fac] {\S\ref{sec:group1}, Thm.~\ref{thm:factor}};
\node[thy, below=1mm of cloud] {\S\ref{sec:group2}};
\node[thy, below=1mm of clus] {\S\ref{sec:group2}};
\node[thy, below=1mm of split] {Thm.~\ref{thm:hm}, Prop.~\ref{prop:bayes-angle}};
\node[thy, below=1mm of out] {\S\ref{sec:labels}, Prop.~\ref{prop:certify}};
\end{tikzpicture}
\caption{The method at a glance. The two grouping steps (thick boxes) are the core of the paper: modes are clustered
onto the curves of the factors within one sample, and factors are clustered across posterior samples. Each box names
the section that describes it and the result that justifies it.}\label{fig:roadmap}
\end{figure}

\paragraph{Contributions}
\begin{enumerate}
  \item \emph{The grouping principle and its limits}: groups are factors, with the exact ambiguity $\deltaP$
  (\cref{thm:factor}); terms are not groups (\cref{prop:terms}); the conditioning of the split (\cref{thm:hm}) and its
  Bayesian form (\cref{prop:bayes-angle}); what finite data can certify (\cref{thm:ident,prop:certify}).
  \item \emph{An algorithm that learns the groups} from noisy, gridded or scattered data, by clustering modes onto
  curves and factors across samples, whose output encodes these limits as intervals and labels.
  \item \emph{Evidence}: on $285$ random equations never used in development, $96.5\%$ right or equivalent and no
  confident wrong answer, against $28.4\%$ and $204$ confident wrong answers for weak SINDy; on a ten-equation
  benchmark in three designs, $30$ of $30$ factor sets.
\end{enumerate}

\paragraph{Relation to prior work}
Most methods that learn a PDE from data regress a time derivative on a library of candidate terms: SINDy and
PDE-FIND~\cite{BruProKut16,RudBruProKut17}, their weak, ensemble and Bayesian variants~\cite{MesBor21,FasKutBru22,HirBarKut22},
and Gaussian-process formulations of linear operators~\cite{RaiPerKar17}. Their unit is the \emph{term}. Dynamic mode
decomposition~\cite{Sch10,SasKut22} works with eigenvalues, as we do, but returns a linear operator on the data, not an
equation, and does not group eigenvalues by the factor they belong to. Matrix pencils~\cite{HuaSar90}, the algebra of
constant-coefficient operators and the Ehrenpreis--Palamodov fundamental principle~\cite{Ehr70,Pal70,Mal56,Hor90,Obe90},
and Dirichlet-process mixtures~\cite{BleJor06} are classical. What is new is the connection: the \emph{factor} of the
symbol, not the term, as the unit of learning; what finite data can certify about it and how stably it can be
separated; and an algorithm that groups twice and reports these limits.

\paragraph{Scope} Everything proved and implemented concerns scalar linear equations with constant coefficients in one
time and one space variable, observed on a rectangle. Equation~\eqref{eq:intro-tpd} with $c=1$, $\nu=0.3$ is the running
example. The proofs of the main results, the complete algorithm with its parameters, further cases, a guide to reading the
output, and the code and benchmark data are in the supplementary materials; the code is public~\cite{Pri26code}.

